%%%%%     PACKS     %%%%%
\documentclass[12pt]{article}
\usepackage[margin=1in,headsep=.60in]{geometry}
\usepackage[utf8]{inputenc}
\usepackage[table, dvipsnames]{xcolor}
\usepackage{array}
\usepackage{amsmath}
\usepackage{booktabs}
\usepackage{mdframed} %For box around text
\usepackage{hyperref} %For hyperlinks in the PDF
\hypersetup{
    colorlinks=true,
    linkcolor=blue,
    filecolor=magenta,
    urlcolor=cyan,
}
\urlstyle{same}
\usepackage{amsthm}
\usepackage{amssymb}
\usepackage{amsfonts}
\usepackage{siunitx}
\usepackage{graphicx}
\usepackage{caption}
\usepackage{pgfplots}
\graphicspath{{Images/}}
\usepackage[colorinlistoftodos]{todonotes}
\usepackage{cleveref}
\usepackage[labelformat=simple]{subcaption}
\usepackage{grffile}
\usepackage{gensymb}
\usepackage{float}
\usepackage[shortlabels]{enumitem}
\usepackage{enumitem}
\setlistdepth{9}
\usepackage{biblatex} %Imports biblatex package
\addbibresource{references.bib}
\usepackage{outlines}
\usepackage{minted}
\usepackage{pdflscape}
\usepackage{everypage}
\usepackage{multirow}
\usepackage{multicol}
\usepackage{afterpage}
\usepackage[labelfont=bf, font=small]{caption}
\usepackage{upgreek}
\usepackage{tikz}
\usepackage{xcolor}
\usepackage{listings}

% style sheet for code
\definecolor{codegreen}{rgb}{0,0.6,0}
\definecolor{codegray}{rgb}{0.5,0.5,0.5}
\definecolor{codepurple}{rgb}{0.58,0,0.82}
\definecolor{backcolour}{rgb}{0.95,0.95,0.92}

\lstdefinestyle{mystyle}{
    backgroundcolor=\color{backcolour},   
    commentstyle=\color{codegreen},
    keywordstyle=\color{magenta},
    numberstyle=\tiny\color{codegray},
    stringstyle=\color{codepurple},
    basicstyle=\ttfamily\footnotesize,
    breakatwhitespace=false,         
    breaklines=true,                 
    captionpos=b,                    
    keepspaces=true,                 
    numbers=left,                    
    numbersep=5pt,                  
    showspaces=false,                
    showstringspaces=false,
    showtabs=false,                  
    tabsize=2
}

\lstset{style=mystyle}

\usepackage{lastpage}
%%%%%     COMMANDS     %%%%%
%%%% Blue box for subsection text (not figures) %%%%
\newenvironment{bluebox}
  {\begin{mdframed}[backgroundcolor=blue!5,linecolor=blue!40,roundcorner=8pt]}
  {\end{mdframed}}

%%%% Simple neutral box for figures %%%%
\newenvironment{figbox}
  {\begin{mdframed}[roundcorner=8pt,shadow=true,shadowsize=4pt,shadowcolor=black!40]}
  {\end{mdframed}}

\begin{document}

\begin{titlepage}
\newcommand{\HRule}{\rule{\linewidth}{0.5mm}}

\center

\textsc{\LARGE Aarhus University}\\[1.5cm]
\textsc{\Large Department Of Mathematics}\\[0.5cm]
\textsc{\large Introduction to Algebra 2}\\[0.5cm]
    

\HRule\\[0.4cm]
    \center 
    {\huge\bfseries  Handin 1}\\[0.4cm] % Title of your document
\HRule\\[1.5cm]

\begin{minipage}{0.4\textwidth}
        \begin{flushleft}
            \large
            \textit{Author}\\   
            Søren M. \textsc{Damsgaard}\\
                 % Your name
        \end{flushleft}
    \end{minipage}
~
    \begin{minipage}{0.4\textwidth}
        \begin{flushright}
            \large
            \textit{Student number}\\
            \textbf{202309814}\\
                 % Studienummer\
            
        \end{flushright}
    \end{minipage}

\vfill\vfill\vfill % Position the date 3/4 down the remaining page
    
    {\large\today}

\vfill\vfill
    \includegraphics[width=0.2\textwidth]{Aarhus_University_seal.png}\\[1cm] % Include a department/university logo - this will require the graphicx package
    

\vfill
\end{titlepage}
%%%%%      CHAPTERS      %%%%%
\hspace{0.02cm}
\section*{Exercise 13.5}
Show Theorem 13.3~\cite{Algebra}.
\subsection*{Show first that $K[X]$ has at least one monic irreducible polynomial.}
First, a quick definition of what we are working with:
\begin{enumerate}
    \item $K$ is a field, see Definition 12.4~\cite{Algebra}.
    \item $K[X]$ is a polynomial ring, see Chapter 13.1~\cite{Algebra}.
    \item Monic is when the leading coefficient of a polynomial is 1, see Chapter 13.1~\cite{Algebra}.
    \item Irreducible is when every factorization of a non-constant polynomial in $K[X]$ only has trivial factorizations (where one factor is a unit), see Section 13.2~\cite{Algebra}.
\end{enumerate}

Let's start by showing that a monic polynomial exists in $K[X]$.\\
Per the definition of a ring 12.1(f)~\cite{Algebra}, there exists a neutral element $1_K \in K$.\\
Now, let there be a polynomial of the first degree, $X^1$
\begin{equation*}
    X^1 \cdot 1 = 1 \cdot X^1 = X^1
\end{equation*} 
Here, it holds that the leading coefficient is $1$, so the polynomial is monic!\\
\newline
Now we must show that the monic polynomial $X^1$ is also irreducible.\\
This is done by showing that all factorizations
\begin{equation} \label{trivial}
    X^1 = f(X)\cdot g(X)
\end{equation}
are trivial. A factorization is called trivial if one of the factors is a unit, definition 12.8~\cite{Algebra}.\\
According to Theorem 13.1(a)~\cite{Algebra}, it holds that the degrees of $f(X)$ and $g(X)$ must sum to the degree of $X^1$:
\begin{equation*}
    \deg f(X) + \deg g(X) = 1 \qquad \text{for } \deg f(X), \deg g(X) \geq 0
\end{equation*}
This means at least one of them must have a degree of $0$ and be a constant. It holds from Theorem 13.1(c)~\cite{Algebra}, a polynomial in $K[X]$ is a unit if and only if it is a constant and a unit in $K$.\\
 $K$ is a field where all non-zero elements are units (Definition 12.4)~\cite{Algebra}, so one of the factors must be a unit in $K[X]$.\\
This makes all factorizations in \eqref{trivial} trivial, making $X^1$ irreducible.\\
This shows that there is AT LEAST ONE monic irreducible polynomial in $K[X]$.

\subsection*{Show that the assumption that a finite number of irreducible monic polynomials exists leads to a contradiction}
This can be shown using the Euclids proof from 5.1 ~\cite{Algebra}.\\
Assume there is a finite number of irreducible monic polynomials, $P_1(X),P_2(X)\hdots P_n(X)$.\\
Construct a new polynomial $F(X)$
\begin{equation*}
    F(X) = P_1(X) \cdot P_2(X) \cdots P_n(X) + 1
\end{equation*} 
This new polynomial is monic according to Theorem 13.1(a) ~\cite{Algebra}, the leading coefficient of $F(X)$ is the product of the coefficient of all its polynomial factors, which we know to be 1, since they're all monic, hence $F(X)$ is monic\\
Now from Theorem 13.2 ~\cite{Algebra}, and because $\deg F(X) > 0$, $F(X)$ then either $F(X)$ is irreducible or one of its factors must be irreducible, in both cases a monic irreducible factor exists, let's call it $q(X)$.\\
\newline
This new factor $q(X)$ must be in the the list $P_1(X),P_2(X)\hdots P_n(X)$ since it's assumed to contain all monic irreducible polynomials, so $q(X) = P_i(X)$ for some $ i \in \{1,2, \hdots ,n\}$.\\
This leads to the following
\begin{align*}
    & P_i(X) \mid F(X) \land P_i(X) \mid P_1(X) \cdot P_2(X) \cdots P_n(X) \\
    & \implies P_i(X) \mid F(X) - (P_1(X) \cdot P_2(X) \cdots P_n(X)) \\
    & \implies P_i(X) \mid 1
\end{align*}
This leads to a contradiction since the polynomials that can divide 1 is the constant polynomials, but since $P_i(X)$ is irreducible it can't be constant, which means the assumption of there being a finite number of monic irreducible polynomials must be false and there must be an infinite amount of monic irreducible polynomials.

\printbibliography
\end{document}